\section{Phạm vi của phương pháp}\label{sec:02-scope}
Phương pháp trong chương này có ba giới hạn cần giữ khi đọc kết quả. Thứ nhất, dữ liệu phân cụm và xếp hạng là bán tổng hợp hoặc tổng hợp, nên kết quả RQ1--RQ3 mô tả hành vi của thuật toán dưới các điều kiện mô phỏng có kiểm soát, không phải hiệu quả trên báo cáo cứu hộ thực tế. Thứ hai, bộ ảnh do nhóm thu thập và gán nhãn, chưa có giấy phép cho từng ảnh, chưa có mẫu chụp thực địa độc lập và chưa đo độ đồng thuận giữa người gán nhãn. Thứ ba, các hằng số của $Q_i$ và điểm ưu tiên là heuristic của nhóm, chưa được chuyên gia cứu hộ xác nhận.

Đối với thông tin khẩn cấp, đề tài sử dụng Logistic Regression từ các lựa chọn có cấu trúc trên biểu mẫu, thay cho phương án phân tích văn bản tiếng Việt tự do dự kiến ban đầu. Đây là surrogate trên năm cờ với nhãn OR tự sinh, không phải xác suất nguy kịch đã được hiệu chuẩn trên dữ liệu thực. Trong phạm vi repository và các ref đã khảo sát, chưa tìm thấy mô hình văn bản tiếng Việt tự do được huấn luyện hoặc đánh giá như cam kết DistilBERT. Nhánh NLP của ứng dụng mới có tiếp nhận mô tả và heuristic từ khóa; chưa chứng minh hoàn thành cam kết nghiên cứu. Pilot văn bản tiếng Anh bổ sung ở Mục~\ref{sec:multimodal-pilot} có nhiệm vụ riêng, không thay thế đánh giá tiếng Việt hoặc nguy cấp.

Cần phân biệt ba mức: tiếp nhận ảnh, mô tả và vị trí; tích hợp các trường bằng công thức, ánh xạ và heuristic; học một mô hình fusion trên dữ liệu ghép cặp. Ứng dụng cứu hộ có hai mức đầu với khoảng trống nối điểm LR vào adapter. Đợt bổ sung đã huấn luyện fusion đặc trưng tuyến tính trên CrisisMMD, chưa tích hợp vào ứng dụng và chưa fine-tuning chung hai encoder. Nguồn tham chiếu dấu hiệu IITT không xác nhận các trọng số điều phối của nhóm. Sửa cách mô tả phạm vi và benchmark bổ trợ không đồng nghĩa đã hoàn thành các cam kết còn thiếu.

Mục~\ref{sec:text-design} và Mục~\ref{sec:multimodal-design} bổ sung kiến trúc, đầu vào/đầu ra và quy trình xác minh cho văn bản và fusion tiếng Việt. Đây là thiết kế nghiên cứu chưa thực hiện, với target đề xuất là có/không có yêu cầu hỗ trợ cứu hộ; chưa phải phân loại mức nguy cấp hoặc chính sách ưu tiên đã xác nhận. Nội dung nghiên cứu được đối chiếu với thuyết minh ở Mục~\ref{sec:proposal-comparison}; giới hạn của các phép thử được trình bày ở Mục~\ref{sec:04-evaluation-limitations}.

Ablation ảnh/văn bản/cả hai đã thực hiện cho nhiệm vụ informative của CrisisMMD, với target và test IDs chung. Đối với nhiệm vụ ưu tiên cứu hộ tiếng Việt, vẫn thiếu report ID ghép cặp và gold độc lập chung. Mask các trường $F/E$ mô phỏng chỉ là phân tích độ nhạy nguồn thông tin có cấu trúc, không thay thế ablation mô hình ảnh--văn bản. Cần vượt các điều kiện dữ liệu và nhãn trước khi khóa protocol cứu hộ, không đổi target để coi benchmark bổ trợ là bằng chứng cùng nhiệm vụ.
